% ---- pg-292 （书页 280） ----

\chapter{多元线性回归模型}
\phantomsection\label{chap:7}

\section{引言}
\phantomsection\label{sec:7.1}

回归分析是一种通过一组预测变量(自变量)来预测一个或多个响应变量(因变量)的统计方法. 它也可用于评估预测变量对响应变量的效果. 不幸的是, 回归这一名称来源于 F. 高尔顿(见文献[15])就这一研究课题发表的第一篇论文的题目, 既没有充分反映这种方法的重要性, 也没有反映这种方法应用的广泛性.

在本章中, 我们首先讨论预测一个响应变量的多重回归\footnote{{\footnotesize 译者注: 这里, ``multiple regression'' 是指有多个自变量(回归变量)的回归模型, 我们译为``多重回归'', 而 ``multivariate regression'' 是指有多个响应变量(因变量)的回归模型, 我们译为``多元回归''. 读者应注意这二者的区别.}}, 然后将这个模型推广到预测若干个因变量的情形. 由于在这方面存在大量文献, 我们的讨论肯定是简明扼要的. (若你有兴趣深入研究回归分析, 可参考以下各种按难度递增顺序排列的著作: Abraham 和 Ledolter(见文献[1]), Bowerman 和 O'Connell(见文献[6]), Neter, Wasserman, Kutner 和 Nachtsheim(见文献[20]), Draper 和 Smith(见文献[13]), Cook 和 Weisberg(见文献[11]), Seber(见文献[23]), 以及 Goldberger(见文献[16]).) 我们对回归分析的简要讨论重点在以下几方面: 模型的假定及其后果, 回归模型的一些不同表述方式, 以及回归方法在看上去非常不同的场合下的一般应用.

\section{经典线性回归模型}
\phantomsection\label{sec:7.2}

设 $z_1,z_2,\cdots,z_r$ 为 $r$ 个预测变量, 设想它们与一个响应变量 $Y$ 有关系. 例如, 对 $r=4$, 我们可能有
\[
\begin{aligned}
Y &= \text{住房的当前市场价值}\\
z_1 &= \text{居住面积(平方英尺)}\\
z_2 &= \text{位置(城市区域的指标)}\\
z_3 &= \text{去年的评估价值}\\
z_4 &= \text{建筑质量(每平方英尺价格)}
\end{aligned}
\]
经典线性回归模型假定 $Y$ 由一个均值和一个随机误差 $\epsilon$ 合成, 其中均值为诸 $z_i$ 的连续函数, 而 $\epsilon$ 则考虑测量误差和其余未被明确考虑在模型中的变量所产生的效应. 在实验中被记录下来的或由研究者设定的预测变量的值被视为固定值. 误差被看成是随机变量(因而响应变量
